\scanpage{60}
\textcircled{18} À tout $B\in(\mathrm{Ens})_6$, ensemble de cardinal 6, associons le graphe $\unsure{\mathscr T|\Gamma|\mathcal T}_1(B)$ suivant
\[
\left\{\begin{aligned}
&\text{ens sommets}=\unsure{\mathfrak P|\mathscr P}_2(B)\\
&\deleted{A,B,}\ u,v\in\mathfrak P_2(B)\quad\text{liés ssi}\quad u\cap v=\varnothing
\end{aligned}\right.
\]

Alors $\mathscr T_1(B)$ est un graphe \unsure{trilétaire|trialitaire|trilitaire|tritétraèdre} donc un $\mathscr T_1$, et le foncteur $B\longrightarrow\mathscr T_1(B)$ est une équivalence de la catégorie $(\mathrm{Ens})_6$ avec la catégorie des $\mathscr T_1$. Le foncteur quasi-inverse associe à $\unsure{\mathscr X|\mathfrak X|\mathcal X}$ l'ens $\underline{B}$ des parties discrètes de card 5 de $\mathscr X$.
Alors tt $x\in\mathscr X$ est contenu dans exactement deux éléments de $B$, et l'application $\mathscr X\mathbin{\unsure{-|\longrightarrow}}\mathfrak P_2(B)$ ainsi obtenue est un isom. du graphe $\mathscr X$ avec le graphe $\mathscr T_1(B)$.
La donnée d'un bitriangle de $\mathscr X$ équivaut à la donnée d'une partition de type $3,3$ de $B$ -- à celle-ci correspond l'ens des \deleted{\unsure{parties}} \inserted{$\in\mathfrak P_2(B)$} qui sont des sections de $B$ sur l'ens \unsure{quot\textsuperscript{t}} à deux éléments.
Tout bitriangle a un $\mathscr T_2$ enveloppe canonique (le foncteur des $(\mathscr X,X)$ d'un $\mathscr T_1$ et d'un $\mathscr T_0$ contenu, vers les $X$, est une équivalence de catégories).

% archive sheet number: 59

